% ==STYLE-BEGIN==
\documentclass[UTF8,11pt]{ctexart}
\usepackage{amsmath,amssymb,amsthm}
\usepackage[a4paper,margin=2.4cm,top=2.3cm,bottom=2.3cm]{geometry}
\usepackage{fancyhdr}
\usepackage{enumitem}
\usepackage{hyperref}
\hypersetup{hidelinks}

% ---- 版式 ----
\linespread{1.15}
\setlength{\parskip}{0.35em}
\setlength{\parindent}{1.6em}
\setlist{nosep,leftmargin=2em}

% ---- 页眉 ----
\pagestyle{fancy}
\fancyhf{}
\fancyhead[L]{\small\coursename}
\fancyhead[R]{\small\thepage}
\renewcommand{\headrulewidth}{0.4pt}

% ---- 定理环境：中文用正体，标题加粗，按小节共用一条流水号 ----
\newtheoremstyle{cnthm}{0.7em}{0.7em}{\normalfont}{}{\bfseries}{}{0.5em}{}
\theoremstyle{cnthm}
\newtheorem{thm}{定理}[section]
\newtheorem{lem}[thm]{引理}
\newtheorem{prop}[thm]{命题}
\newtheorem{cor}[thm]{推论}
\newtheorem{defn}[thm]{定义}
\newtheorem{rem}[thm]{注}
\newtheorem{exa}[thm]{例}
\newtheorem{conj}[thm]{猜想}
\renewcommand{\proofname}{\bfseries 证明}
\providecommand{\coursename}{}
% ==STYLE-END==
\renewcommand{\coursename}{解析几何\,·\,0930}
\title{几何与代数 0930}
\date{}
\begin{document}
\maketitle

\section{\ 二次曲面}

\textbf{二次曲面}：由一个三元二次方程确定的空间图形。

三元二次方程
\[
a_{11}x^2+a_{22}y^2+a_{33}z^2
+2a_{12}xy+2a_{13}xz+2a_{23}yz
+2a_1x+2a_2y+2a_3z+c=0
\]
可以写成如下的矩阵形式
\[
(x,y,z,1)
\begin{pmatrix}
a_{11} & a_{12} & a_{13} & a_1\\
a_{12} & a_{22} & a_{23} & a_2\\
a_{13} & a_{23} & a_{33} & a_3\\
a_1 & a_2 & a_3 & c
\end{pmatrix}
\begin{pmatrix}
x\\y\\z\\1
\end{pmatrix}
=0 .
\]

\textbf{Q：}三元二次方程表示什么图形？图形有几类？

平面上的二次曲线有几类？
\[
a_{11}x^2+a_{22}y^2+2a_{12}xy+2a_1x+2a_2y+b=0 .
\]
曲面方程 $F(x,y,z)=0$；曲线方程可以看作两个曲面相交：
$F(x,y,z)=0$，$G(x,y,z)=0$。

\section{\ 球面}

\textbf{球面}：到定点的距离等于定长的点的集合。

设定点 $Q_0(x_0,y_0,z_0)$，定长为 $r$，任取点 $P(x,y,z)$，
则 $|PQ_0|=r$，有
\[
\Longrightarrow\ (x-x_0)^2+(y-y_0)^2+(z-z_0)^2=r^2
\]
\[
\Longleftrightarrow\
x^2+y^2+z^2-2x_0x-2y_0y-2z_0z
+x_0^2+y_0^2+z_0^2-r^2=0 .
\]
该方程有如下特点：

\begin{enumerate}
\item 没有交叉项；
\item 平方项系数相同。
\end{enumerate}

反过来，若三元二次方程满足上两条件，即方程为
\[
x^2+y^2+z^2+2a_1x+2a_2y+2a_3z+b=0,
\]
可化为
\[
(x+a_1)^2+(y+a_2)^2+(z+a_3)^2
=a_1^2+a_2^2+a_3^2-b .
\]
于是，该方程表示球心在 $(-a_1,-a_2,-a_3)$、半径为
$\sqrt{a_1^2+a_2^2+a_3^2-b}$ 的球面。

\subsubsection*{e.g.1}

求过 $(A,0,0)$、$(0,B,0)$、$(0,0,C)$ 的球面方程。

设过三点的球面方程为
\[
x^2+y^2+z^2+2a_1x+2a_2y+2a_3z+b=0 .
\]
代入三点坐标有
\[
\begin{cases}
A^2+2a_1A+b=0,\\
B^2+2a_2B+b=0,\\
C^2+2a_3C+b=0 .
\end{cases}
\]
令 $b=0$ 得
\[
a_1=-\frac A2,\qquad a_2=-\frac B2,\qquad a_3=-\frac C2 .
\]
所以方程为
\[
x^2+y^2+z^2-Ax-By-Cz=0 .
\]

\section{\ 柱面}

\textbf{柱面}：由一族平行直线所组成的图形。平面是柱面。

这些直线称为柱面的\textbf{母线}；与每一条母线相交的曲线叫
\textbf{准线}。柱面是由一条母线沿着准线平行移动得到的图形。
由此看来，确定一个柱面，需要准线 $+$ 母线方向。

\textbf{圆柱面}：到定直线的距离为定长的点的集合，
且准线是圆、母线方向与圆垂直的柱面。

取 $[O;X,Y,Z]$，s.t. $z$ 轴是母线方向，柱面 $S$ 的方程
$F(x,y,z)=0$，$\forall P(x,y,z)\in S$，则
$\forall t\in\mathbb R$，$(x,y,tz)\in S$，即 $F(x,y,tz)=0$
$\forall t\in\mathbb R$ 都成立，于是柱面 $S$ 的方程为 $F(x,y)=0$。

反之，如果 $F(x,y,z)=0$ 不含 $z$，则该方程 $F(x,y)=0$
表示什么？若 $P(x,y,z)\in$ 图形 $S$，则
$\forall t\in\mathbb R$，$(x,y,tz)\in S$，说明过 $P$ 点且平行于
$z$ 轴的直线属于 $S$，于是 $S$ 是母线方向为 $z$ 轴的柱面。

柱面的准线方程
\[
\Gamma=
\begin{cases}
F(x,y)=0,\\
z=0 .
\end{cases}
\]

\textbf{圆柱面方程}\quad 取 $[O;X,Y,Z]$，s.t. $O$ 是准线的圆心，
$z$ 轴是母线方向。已知准线 $\Gamma$，
$\forall P(x,y,z)$，过 $P$ 点与 $z$ 轴平行的直线交准线于
$(x,y,0)$，于是 $x^2+y^2=r^2$。

还有很多柱面，如椭圆柱面
$\dfrac{x^2}{a^2}+\dfrac{y^2}{b^2}=1$，双曲柱面
$\dfrac{x^2}{a^2}-\dfrac{y^2}{b^2}=1$，抛物柱面 $y^2=2px$。

\section{\ 锥面}

\textbf{锥面}：由一族交于同一点的直线组成的图形。平面是锥面。

该点称为\textbf{顶点}；这些直线叫\textbf{母线}；
与所有母线相交（顶点除外）的曲线叫做\textbf{准线}。
由此看来，确定一个锥面，需要顶点 $+$ 准线。

\textbf{圆锥面}：准线是圆，顶点在过准线的圆心且垂直于圆的直线上。

取 $[O;X,Y,Z]$，s.t. $O$ 是锥面 $S$ 的顶点，$S$ 的方程
$F(x,y,z)=0$。$\forall P(x,y,z)\in S$，则直线 $OP$ 也在锥面
$S$ 上，即 $\forall t\in\mathbb R$，$(tx,ty,tz)\in S$，于是
\[
F(tx,ty,tz)=0 .
\]
锥面方程特点：若 $F(x,y,z)=0$，则
$F(tx,ty,tz)=0$，$\forall t\in\mathbb R$。

反之，若方程满足以上条件，它表示什么图形？
$\forall P(x,y,z)\in S$，由于 $F(x,y,z)=0
\Longrightarrow F(tx,ty,tz)=0$，说明
$(tx,ty,tz)\in S$，$\forall t\in\mathbb R$，于是直线
$OP\subset S$，可知 $S$ 是顶点在原点的锥面。

\end{document}
